%
% Front matter
%

\begin{frontmatter}

\title{
  %T2K \TotalPotNu-POT Neutrino  \TotalPotNuBar-POT Anti-Neutrino 3-Flavour Joint Analysis
  T2K Neutrino and Anti-Neutrino 3-Flavour Joint Analysis of  \\
  Run 1-6 (\TotalPotNu-POT $\nu$  \TotalPotNuBar-POT $\bar{\nu}$), \\
  Run 1-7b (\TotalPotNuRunOneSevenB-POT $\nu$  \TotalPotNuBarRunOneSevenB-POT $\bar{\nu}$) and \\
  Run 1-7c (\TotalPotNuRunOneSevenC-POT $\nu$  \TotalPotNuBarRunOneSevenC-POT $\bar{\nu}$)
  data sets
}

\vspace*{0.7cm}

\author[RAL,LIV]  {Costas Andreopoulos}
\author[LIV]  {Christopher Barry}
\author[LANCS]  {Thomas Dealtry}
\author[RAL,LIV]  {Steve Dennis\\}
%\author[IFIC]     {Lorena Escudero S\'anchez}
%\author[WAR]{Nick Grant}
\author[GEN]     {Davide Sgalaberna}
\author[RAL,OXF]  {Raj Shah}
%\author[RAL,OXF]  {and Alfons Weber}

\address[RAL]   {Science \& Technology Facilities Council (STFC), Rutherford Appleton Laboratory,\\ Harwell Oxford Campus, OX11 0QX, UK}
\address[LIV]   {University of Liverpool, Physics Department,\\ Liverpool, L69 7ZE, UK} 
%\address[ETHZ]  {ETH Zurich, Institute for Particle Physics,\\ 8092 Zurich, Switzerland}
\address[OXF]   {University of Oxford, Subdepartment of Particle Physics,\\ Oxford, OX1 3RH, UK}
%\address[WAR]   {University of Warwick, Physics Department,\\ Coventry, CV4 7AL, UK}
%\address[IFIC]  {CSIC and University of Valencia, Instituto de Fisica Corpuscular, \\ E-46071 Valencia, Spain}
\address[LANCS] {Lancaster University, Physics Department,\\ Lancaster, LA1 4YB, UK}
\address[GEN]  {University of Geneva, Departement de physique nucleaire et corpusculaire,\\ 1211 Geneve, Switzerland}

%
% Abstract
%

\begin{abstract}

We report the results of the
3-flavour $\nu/\bar{\nu}$ joint analysis, performed by the {\em VALOR} group,
on the combined Run 1+6 dataset as well as on the combined Run 1+7 dataset, 
using data taken in both neutrino and antineutrino beam mode.
The goal of this analysis is 
to search for CP violation by setting limits on $\delta_{CP}$
and measure the oscillation parameters
$\sin^2 \theta_{13}$, $\sin^2 \theta_{23}$, $\Delta m^2_{32}$ ($\Delta m^2_{13}$) 
($\Delta m^2_{32}$ for normal hierarchy and $\Delta m^2_{13}$ for inverted hierarchy).
Though T2K has limited sensitivity to the mass hierarchy given the relatively short baseline,
fits under both hypotheses are performed to show which is preferred by data.
Furthermore the search for $\bar{\nu}_e$ appearance and 
the measurement of $\bar{\nu}_{\mu}$ disappearance are shown as well.

This analysis adapts the work presented in Refs. 
\cite{OurRun1DisappAnalysis},
\cite{OurRun12DisappAnalysis}, 
\cite{OurRun12DisappAnalysis3Flavour},
\cite{OurRun123DisappAnalysis3Flavour},
\cite{OurRun1234aDisappAnalysis3Flavour},
\cite{OurRun1234DisappAnalysis3Flavour},
\cite{OurRun123JointOscillationAnalysis},
\cite{OurRun1234JointOscillationAnalysis}
\cite{OurRun56AntiNumuDisappearanceAnalysis3Flavour} 
and
\cite{OurRun56AntiNueBarAppearanceAnalysis3Flavour}
to this dataset.

% Run 1-6

The Run 1-6 dataset corresponds to an integrated J-PARC neutrino beam
exposure of \TotalPotNu{} POT in neutrino mode and
\TotalPotNuBar{} POT in antineutrino mode.
This analysis predicts 
\NskOscExpectationNumu{} $\pm$ \NskOscExpectationNumuError{} (syst) for neutrino $\mu$-like,
\NskOscExpectationNumuBar{} $\pm$ \NskOscExpectationNumuBarError{} (syst) for antineutrino $\mu$-like,
\NskOscExpectationNue{} $\pm$ \NskOscExpectationNueError{} (syst) for neutrino $e$-like,
and
\NskOscExpectationNueBar{} $\pm$ \NskOscExpectationNueBarError{} (syst) for antineutrino $e$-like
single-ring events in Super-K,
given the oscillation parameters of the Asimov data set A in \cref{tab:osc_params_joint_asimov},
corresponding to $\sin^2 \theta_{23} = 0.528$, $\delta_{CP} = -1.601$ and normal hierarchy, 
and observes respectively 
\NskObservedNumu{} for neutrino $\mu$-like,
\NskObservedNumuBar{}  for antineutrino $\mu$-like,
\NskObservedNue{}  for neutrino $e$-like,
\NskObservedNueBar{}  for antineutrino $e$-like.

% Run 1-7b

The Run 1-7b dataset corresponds to an integrated J-PARC neutrino beam
exposure of \TotalPotNuRunOneSevenB{} POT in neutrino mode and
\TotalPotNuBarRunOneSevenB{} POT in antineutrino mode.

This analysis predicts 
\NskOscExpectationNumuRunOneSevenB{} $\pm$ \NskOscExpectationNumuErrorRunOneSevenB{} (syst) for neutrino $\mu$-like,
\NskOscExpectationNumuBarRunOneSevenB{} $\pm$ \NskOscExpectationNumuBarErrorRunOneSevenB{} (syst) for antineutrino $\mu$-like,
\NskOscExpectationNueRunOneSevenB{} $\pm$ \NskOscExpectationNueErrorRunOneSevenB{} (syst) for neutrino $e$-like,
and
\NskOscExpectationNueBarRunOneSevenB{} $\pm$ \NskOscExpectationNueBarErrorRunOneSevenB{} (syst) for antineutrino $e$-like
single-ring events in Super-K,
given the oscillation parameters of the Asimov data set A in \cref{tab:osc_params_joint_asimov},
corresponding to $\sin^2 \theta_{23} = 0.528$, $\delta_{CP} = -1.601$ and normal hierarchy, 
and observes respectively 
\NskObservedNumuRunOneSevenB{} for neutrino $\mu$-like,
\NskObservedNumuBarRunOneSevenB{}  for antineutrino $\mu$-like,
\NskObservedNueRunOneSevenB{}  for neutrino $e$-like,
\NskObservedNueBarRunOneSevenB{}  for antineutrino $e$-like.

% Run 1-7c

The Run 1-7c dataset corresponds to an integrated J-PARC neutrino beam
exposure of \TotalPotNuRunOneSevenC{} POT in neutrino mode and
\TotalPotNuBarRunOneSevenC{} POT in antineutrino mode.

This analysis predicts 
\NskOscExpectationNumuRunOneSevenC{} $\pm$ \NskOscExpectationNumuErrorRunOneSevenC{} (syst) for neutrino $\mu$-like,
\NskOscExpectationNumuBarRunOneSevenC{} $\pm$ \NskOscExpectationNumuBarErrorRunOneSevenC{} (syst) for antineutrino $\mu$-like,
\NskOscExpectationNueRunOneSevenC{} $\pm$ \NskOscExpectationNueErrorRunOneSevenC{} (syst) for neutrino $e$-like,
and
\NskOscExpectationNueBarRunOneSevenC{} $\pm$ \NskOscExpectationNueBarErrorRunOneSevenC{} (syst) for antineutrino $e$-like
single-ring events in Super-K,
given the oscillation parameters of the Asimov data set A in \cref{tab:osc_params_joint_asimov},
corresponding to $\sin^2 \theta_{23} = 0.528$, $\delta_{CP} = -1.601$ and normal hierarchy, 
and observes respectively 
\NskObservedNumuRunOneSevenC{} for neutrino $\mu$-like,
\NskObservedNumuBarRunOneSevenC{}  for antineutrino $\mu$-like,
\NskObservedNueRunOneSevenC{}  for neutrino $e$-like,
\NskObservedNueBarRunOneSevenC{}  for antineutrino $e$-like.



A joint analysis of all the four selected samples is performed 
in a framework of 3-flavour oscillations including matter effects in constant-density matter.

For $\mu$-like samples the templates are built as a function of the reconstructed neutrino energy,
while for $e$-like the analysis templates are built as a function of the reconstructed neutrino energy and lepton angle,
in order to improve the separation of different flavour/reaction components
and increase the sensitivity to $\delta_{CP}$.
The observed single-ring events spectra are compared to the expected spectra 
in order to infer all the 3-flavour oscillation parameters.

As in \cite{OurRun56AntiNueBarAppearanceAnalysis3Flavour},
all the systematic parameters are marginalised (see \cref{sec:oscfit})
%\cite{CousinsHighland} 
in the fit and 
confidence intervals are quoted for all the oscillation parameters.
Since the most important result of this analysis release is the search for CP violation,
the Feldman-Cousins method \cite{FeldmanCousins} is performed only for the simultaneous fit of $\delta_{CP}$ and mass hierarchy.


\end{abstract}

\end{frontmatter}
